\documentclass{../nndlcourse}


\begin{document}


\title{线性分类模型}
\author{数据科学与大数据技术专业}
\date{第 3 讲  }
\maketitle


% ===================== 分类问题与线性模型 =====================
\section{分类问题与线性模型}
\begin{enumerate}
    \item \textbf{分类任务的基本特征}
    \begin{itemize}
        \item 分类任务的目标是将输入映射到离散标签空间, 输出属于有限个类别之一.
        \item 与回归相比, 分类问题更关心决策边界和类别可分性; 在实际应用中, 决策边界往往依赖非线性特征表示.
        \item 典型应用包括图像识别、文本分类、医疗诊断和风险预警等.
    \end{itemize}
    \item \textbf{线性分类模型的判别方式}
    \begin{itemize}
        \item 线性分类器通常先计算线性得分
        \begin{equation*}
            f_c(\boldsymbol{x}) = \boldsymbol{w}_c^\top \boldsymbol{x} + b_c,
        \end{equation*}
        再根据得分输出类别.
        \item 在二分类中, 一个超平面即可将特征空间划分为两个半空间.
        \item 在多分类中, 常见策略包括:
        \begin{itemize}
            \item \textbf{一对多(One-vs-Rest, OvR)}: 为每个类别训练一个二分类器, 预测时选择得分最高的类别.
            \item \textbf{一对一(One-vs-One, OvO)}: 为任意两类分别训练一个二分类器, 共需 $\frac{C(C-1)}{2}$ 个分类器, 预测时通过投票决定类别.
            \item \textbf{直接 $\arg\max$ 判别}: 同时为 $C$ 个类别建立得分函数, 将 $\arg\max_c f_c(\boldsymbol{x})$ 作为预测结果.
        \end{itemize}
    \end{itemize}
    \item \textbf{分类评价指标}
    \begin{itemize}
        \item 给定测试集 $\{(\boldsymbol{x}^{(n)}, y^{(n)})\}_{n=1}^N$ 和预测结果 $\hat{y}^{(n)}$, 准确率与错误率分别为
        \begin{equation*}
            \mathrm{Accuracy} = \frac{1}{N} \sum_{n=1}^N \mathbb{I}\bigl(y^{(n)} = \hat{y}^{(n)}\bigr),
            \qquad
            \mathrm{ErrorRate} = 1 - \mathrm{Accuracy}.
        \end{equation*}
        \item 在二分类中, 混淆矩阵刻画了正负类上的局部表现:
        \begin{center}
        \begin{tabular}{|c|c|c|}
        \hline
         & 预测为正类 ($\hat{y} = 1$) & 预测为负类 ($\hat{y} = 0$) \\
        \hline
        真实为正类 ($y = 1$) & 真正例 (True Positive, TP) & 假负例 (False Negative, FN) \\
        \hline
        真实为负类 ($y = 0$) & 假正例 (False Positive, FP) & 真负例 (True Negative, TN) \\
        \hline
        \end{tabular}
        \end{center}
        \item 基于混淆矩阵, 常用指标包括:
        \begin{itemize}
            \item \textbf{精确率(Precision)}:
            \begin{equation*}
                \mathrm{Precision} = \frac{TP}{TP + FP}.
            \end{equation*}
            它衡量所有“预测为正”的样本中有多少确为正类.
            \item \textbf{召回率(Recall)} 或真正例率(True Positive Rate, TPR):
            \begin{equation*}
                \mathrm{Recall} = \mathrm{TPR} = \frac{TP}{TP + FN}.
            \end{equation*}
            它衡量所有真实正类样本中有多少被模型成功识别.
            \item \textbf{$F_1$ 值}: 精确率与召回率的调和平均,
            \begin{equation*}
                F_1 = \frac{2 \times \mathrm{Precision} \times \mathrm{Recall}}{\mathrm{Precision} + \mathrm{Recall}}
                = \frac{2TP}{2TP + FP + FN}.
            \end{equation*}
            \item \textbf{$F_\beta$ 值}: 用于调节对精确率和召回率的偏好,
            \begin{equation*}
                F_\beta = \frac{(1 + \beta^2)\mathrm{Precision}\mathrm{Recall}}{\beta^2 \mathrm{Precision} + \mathrm{Recall}},
                \qquad \beta > 0.
            \end{equation*}
            其中 $\beta > 1$ 时更偏重召回率, $0 < \beta < 1$ 时更偏重精确率.
        \end{itemize}
        \item \textbf{ROC 曲线与 AUC}: 对输出为概率或打分的二分类器, 可通过调整阈值来考察不同工作点下的性能.
        \begin{itemize}
            \item 在阈值 $t$ 下,
            \begin{equation*}
                \mathrm{TPR}(t) = \mathbb{P}\bigl(s(\boldsymbol{x}) \ge t \mid y = 1\bigr),
                \qquad
                \mathrm{FPR}(t) = \mathbb{P}\bigl(s(\boldsymbol{x}) \ge t \mid y = 0\bigr).
            \end{equation*}
            以 $\mathrm{FPR}$ 为横轴、$\mathrm{TPR}$ 为纵轴即可得到 ROC 曲线.
            \item ROC 曲线越靠近左上角 $(0,1)$, 模型区分能力通常越强; 若曲线接近对角线 $y = x$, 则模型与随机猜测接近.
            \item 若正负样本的打分分布相同, 则对任意阈值都有 $\mathrm{TPR}(t) = \mathrm{FPR}(t)$, 因而 ROC 曲线退化为对角线.
            \item AUC 是 ROC 曲线下的面积,
            \begin{equation*}
                \mathrm{AUC} = \int_0^1 \mathrm{TPR}(u)\, \mathrm{d}u,
                \qquad u = \mathrm{FPR}.
            \end{equation*}
            它衡量模型在所有阈值下的整体排序能力.
            \item 更准确地说, AUC 还可解释为
            \begin{equation*}
                \mathrm{AUC}
                = \mathbb{P}\bigl(s(x^+) > s(x^-)\bigr)
                + \frac{1}{2}\mathbb{P}\bigl(s(x^+) = s(x^-)\bigr),
            \end{equation*}
            即随机抽取一个正样本和一个负样本时, 正样本得分不低于负样本的概率.
        \end{itemize}
    \end{itemize}
    \item \textbf{典型线性分类器}
    \begin{itemize}
        \item \textbf{Logistic 回归}: 用于二分类概率建模.
        \item \textbf{Softmax 回归}: Logistic 回归在多分类问题上的自然推广.
        \item \textbf{感知器}: 基于错误驱动的在线学习算法.
        \item \textbf{支持向量机(SVM)}: 通过最大化间隔提升分类鲁棒性.
    \end{itemize}
\end{enumerate}

\textbf{问题思考}:
\begin{itemize}
    \item 线性模型如何实现多分类?
    \item 感知器与 Logistic 回归的主要区别是什么?
\end{itemize}

% ===================== 信息论与损失函数 =====================
\section{信息论与损失函数}
\begin{enumerate}
    \item \textbf{从极大似然到交叉熵}
    \begin{itemize}
        \item 在分类问题中, 模型通常输出条件概率分布 $P(y \mid \boldsymbol{x}; \boldsymbol{w})$, 而不是直接输出类别标签.
        \item 给定训练集 $\{(\boldsymbol{x}^{(i)}, y^{(i)})\}_{i=1}^n$, 极大似然估计的目标是
        \begin{equation*}
            \max_{\boldsymbol{w}} \prod_{i=1}^n P(y^{(i)} \mid \boldsymbol{x}^{(i)}; \boldsymbol{w}),
        \end{equation*}
        等价地, 也可最大化对数似然
        \begin{equation*}
            \max_{\boldsymbol{w}} \sum_{i=1}^n \log P(y^{(i)} \mid \boldsymbol{x}^{(i)}; \boldsymbol{w}).
        \end{equation*}
        \item 因此, 训练分类模型通常等价于最小化负对数似然(Negative Log-Likelihood, NLL), 这正是交叉熵损失的来源.
    \end{itemize}
    \item \textbf{二分类中的交叉熵}
    \begin{itemize}
        \item 设 $y \in \{0, 1\}$, 并记模型输出的正类概率为
        \begin{equation*}
            p(\boldsymbol{x}) = P(y = 1 \mid \boldsymbol{x}; \boldsymbol{w}).
        \end{equation*}
        则条件概率可统一写为
        \begin{equation*}
            P(y \mid \boldsymbol{x}; \boldsymbol{w}) = p(\boldsymbol{x})^y \bigl(1 - p(\boldsymbol{x})\bigr)^{1-y}.
        \end{equation*}
        \item 对独立同分布样本取负对数似然, 得到经验风险
        \begin{equation*}
            \mathcal{L}(\boldsymbol{w})
            = -\sum_{i=1}^n \log P(y^{(i)} \mid \boldsymbol{x}^{(i)}; \boldsymbol{w})
            = -\sum_{i=1}^n \Bigl[
                y^{(i)} \log p(\boldsymbol{x}^{(i)})
                + \bigl(1 - y^{(i)}\bigr)\log\bigl(1 - p(\boldsymbol{x}^{(i)})\bigr)
            \Bigr].
        \end{equation*}
        \item 若将真实标签看作分布 $[y, 1-y]$, 将模型输出看作分布 $[p(\boldsymbol{x}), 1-p(\boldsymbol{x})]$, 则上式恰好等于这两个分布的交叉熵. 因此, 在分类问题中最小化 NLL 与最小化交叉熵是一回事.
    \end{itemize}
    \item \textbf{信息熵, 交叉熵与 KL 散度}
    \begin{itemize}
        \item 对离散分布 $p(x)$, 信息熵定义为
        \begin{equation*}
            H(p) = -\sum_x p(x) \log p(x).
        \end{equation*}
        \item 若用预测分布 $q(x)$ 去近似真实分布 $p(x)$, 则交叉熵定义为
        \begin{equation*}
            H(p, q) = -\sum_x p(x) \log q(x).
        \end{equation*}
        \item KL 散度定义为
        \begin{equation*}
            \operatorname{KL}(p \| q) = \sum_x p(x) \log \frac{p(x)}{q(x)} = H(p, q) - H(p).
        \end{equation*}
        因为 $H(p)$ 与模型参数无关, 所以最小化交叉熵与最小化 $\operatorname{KL}(p \| q)$ 等价.
        \item \textbf{KL 散度的非负性}: 由不等式
        \begin{equation*}
            \log t \le t - 1, \qquad \forall t > 0,
        \end{equation*}
        令 $t = \frac{q(x)}{p(x)}$, 则对所有满足 $p(x) > 0$ 的 $x$ 有
        \begin{equation*}
            \log \frac{q(x)}{p(x)} \le \frac{q(x)}{p(x)} - 1.
        \end{equation*}
        两边乘以 $p(x)$ 并对 $x$ 求和, 得到
        \begin{equation*}
            \sum_x p(x) \log \frac{q(x)}{p(x)}
            \le
            \sum_x \bigl(q(x) - p(x)\bigr)
            = 0.
        \end{equation*}
        因而
        \begin{equation*}
            \operatorname{KL}(p \| q)
            = \sum_x p(x) \log \frac{p(x)}{q(x)}
            \ge 0.
        \end{equation*}
        等号成立当且仅当 $p(x) = q(x)$ 在所有 $p(x) > 0$ 的位置上成立.
    \end{itemize}
\end{enumerate}

\textbf{问题思考}:
\begin{itemize}
    \item 为什么交叉熵损失适合分类问题?
    \item KL 散度与交叉熵的关系是什么?
\end{itemize}

% ===================== Logistic回归(二分类) =====================
\section{Logistic 回归(二分类)}
\begin{enumerate}
    \item \textbf{模型形式与概率解释}
    \begin{itemize}
        \item 为了书写简洁, 以下默认将偏置项并入特征向量. 于是模型先计算线性得分 $\boldsymbol{w}^\top \boldsymbol{x}$, 再通过 Sigmoid 函数输出正类概率.
        \item Logistic 函数定义为
        \begin{equation*}
            \sigma(z) = \frac{1}{1 + \exp(-z)}.
        \end{equation*}
        因而
        \begin{equation*}
            p(y = 1 \mid \boldsymbol{x}; \boldsymbol{w}) = \hat{y} = \sigma(\boldsymbol{w}^\top \boldsymbol{x}).
        \end{equation*}
        \item 预测时, 常将 $\hat{y}$ 与阈值 $0.5$ 比较, 或直接输出概率作为置信度.
    \end{itemize}
    \item \textbf{损失函数与经验风险}
    \begin{itemize}
        \item Logistic 回归通常采用二元交叉熵作为损失函数.
        \item 在包含 $N$ 个样本的训练集上, 平均经验风险写为
        \begin{equation*}
            \mathcal{R}(\boldsymbol{w}) = -\frac{1}{N} \sum_{n=1}^N \left[
                y^{(n)} \log \hat{y}^{(n)} + \bigl(1 - y^{(n)}\bigr) \log \bigl(1 - \hat{y}^{(n)}\bigr)
            \right].
        \end{equation*}
    \end{itemize}
    \item \textbf{梯度计算与参数更新}
    \begin{itemize}
        \item 由 Sigmoid 函数的导数
        \begin{equation*}
            \sigma'(z) = \sigma(z)\bigl(1 - \sigma(z)\bigr),
        \end{equation*}
        可得经验风险关于参数的梯度为
        \begin{align*}
            \frac{\partial \mathcal{R}(\boldsymbol{w})}{\partial \boldsymbol{w}}
            &= -\frac{1}{N} \sum_{n=1}^N \left(
                \frac{y^{(n)}}{\hat{y}^{(n)}} - \frac{1 - y^{(n)}}{1 - \hat{y}^{(n)}}
            \right) \hat{y}^{(n)}\bigl(1 - \hat{y}^{(n)}\bigr) \boldsymbol{x}^{(n)} \\
            &= \frac{1}{N} \sum_{n=1}^N \bigl(\hat{y}^{(n)} - y^{(n)}\bigr) \boldsymbol{x}^{(n)}.
        \end{align*}
        \item 若采用学习率为 $\alpha$ 的梯度下降法, 则参数更新为
        \begin{equation*}
            \boldsymbol{w}_{t+1}
            \leftarrow
            \boldsymbol{w}_t - \alpha \frac{1}{N} \sum_{n=1}^N \bigl(\hat{y}_{\boldsymbol{w}_t}^{(n)} - y^{(n)}\bigr) \boldsymbol{x}^{(n)}.
        \end{equation*}
    \end{itemize}
    \item \textbf{训练特性与停止策略}
    \begin{itemize}
        \item Logistic 回归会持续根据“预测概率与真实标签之间的残差”调整参数, 即使样本已经被分对, 只要模型仍不够自信, 参数仍会继续更新.
        \item 在线性可分且不加正则化时, Logistic 损失往往没有有限的最优参数, 优化过程会倾向于不断增大 $\|\boldsymbol{w}\|$ 来把概率推向 $0$ 或 $1$.
        \item 因此, 实际训练中常结合 $L_2$ 正则化、早停策略或验证集监控, 以改善泛化能力与数值稳定性.
    \end{itemize}
\end{enumerate}

\textbf{问题思考}:
\begin{itemize}
    \item 为什么 Logistic 回归输出概率?
    \item 交叉熵损失如何影响模型训练?
\end{itemize}

% ===================== Softmax回归(多分类) =====================
\section{Softmax 回归(多分类)}
\begin{enumerate}
    \item \textbf{从二分类到多分类}
    \begin{itemize}
        \item Softmax 回归是 Logistic 回归在多分类问题上的自然推广.
        \item 对每个类别 $c \in \{1, \dots, C\}$, 模型先计算线性得分
        \begin{equation*}
            z_c = \boldsymbol{w}_c^\top \boldsymbol{x} + b_c,
        \end{equation*}
        再通过归一化映射得到类别概率.
        \item 预测时通常选择最大概率对应的类别, 即 $\arg\max_c \hat{y}_c$.
    \end{itemize}
    \item \textbf{Softmax 概率映射}
    \begin{itemize}
        \item 对得分向量 $\boldsymbol{z} = (z_1, \dots, z_C)$, Softmax 定义为
        \begin{equation*}
            \hat{y}_c = \operatorname{softmax}(z_c) = \frac{\exp(z_c)}{\sum_{j=1}^C \exp(z_j)}, \qquad c = 1, \dots, C.
        \end{equation*}
        \item 该映射保证每个 $\hat{y}_c > 0$, 且 $\sum_{c=1}^C \hat{y}_c = 1$, 因而输出可以解释为类别条件概率分布.
    \end{itemize}
    \item \textbf{损失函数与模型优化}
    \begin{itemize}
        \item 若第 $n$ 个样本的真实标签采用 one-hot 向量 $\boldsymbol{y}^{(n)} = (y_1^{(n)}, \dots, y_C^{(n)})$, 模型预测为 $\hat{\boldsymbol{y}}^{(n)} = (\hat{y}_1^{(n)}, \dots, \hat{y}_C^{(n)})$, 则多元交叉熵经验风险可写为
        \begin{equation*}
            \mathcal{R} = -\frac{1}{N} \sum_{n=1}^N \sum_{c=1}^C y_c^{(n)} \log \hat{y}_c^{(n)}.
        \end{equation*}
        \item 当 Softmax 与交叉熵结合使用时, 对每个样本的得分分量 $z_c^{(n)}$ 有
        \begin{equation*}
            \frac{\partial \mathcal{R}}{\partial z_c^{(n)}} = \hat{y}_c^{(n)} - y_c^{(n)},
        \end{equation*}
        这一梯度形式非常简洁, 因而适合用梯度下降及其变种进行优化.
    \end{itemize}
\end{enumerate}

\textbf{问题思考}:
\begin{itemize}
    \item Softmax 回归与 Logistic 回归的联系与区别是什么?
    \item 多分类问题中, Softmax 函数的核心作用是什么?
\end{itemize}

\subsection{补充: 为什么使用 Softmax}
\begin{enumerate}
    \item \textbf{最大熵视角}
    \begin{itemize}
        \item 设得分向量 $x = (x_1, \dots, x_C) \in \mathbb{R}^C$, 希望构造概率分布 $p = (p_1, \dots, p_C)$, 满足 $p_i > 0$, $\sum_{i=1}^C p_i = 1$, 并满足期望约束
        \begin{equation*}
            \sum_{i=1}^C p_i x_i = \mu.
        \end{equation*}
        \item 在仅保留这些已知信息的前提下, 最大熵原则要求求解
        \begin{equation*}
            \max_p \left\{- \sum_{i=1}^C p_i \log p_i\right\}
            \quad \text{s.t.} \quad
            \sum_{i=1}^C p_i = 1, \qquad
            \sum_{i=1}^C p_i x_i = \mu.
        \end{equation*}
        \item 引入拉格朗日乘子后, 可得解的形式为
        \begin{equation*}
            p_i = \frac{\exp(\lambda x_i)}{\sum_{j=1}^C \exp(\lambda x_j)}.
        \end{equation*}
        当 $\lambda = 1$ 时, 就得到标准的 Softmax 形式.
    \end{itemize}
    \item \textbf{平移不变性与数值稳定性}
    \begin{itemize}
        \item 对任意常数 $c \in \mathbb{R}$, 有
        \begin{equation*}
            \operatorname{softmax}(x_i + c)
            = \frac{\exp(x_i + c)}{\sum_{j=1}^C \exp(x_j + c)}
            = \frac{\exp(x_i)}{\sum_{j=1}^C \exp(x_j)}.
        \end{equation*}
        因而
        \begin{equation*}
            \operatorname{softmax}(\boldsymbol{x} + c\mathbf{1}) = \operatorname{softmax}(\boldsymbol{x}).
        \end{equation*}
        \item 这表明 Softmax 只依赖各类别得分之间的相对差异. 实际计算中, 常先减去 $\max_i x_i$ 来避免指数溢出.
    \end{itemize}
    \item \textbf{从 Logistic 到 Softmax}
    \begin{itemize}
        \item 在二分类中, Logistic 回归通过 Sigmoid 函数把一个实数得分映射为正类概率.
        \item 在多分类中, 我们需要把 $C$ 个类别的得分同时映射为和为 $1$ 的概率分布. 对每个得分取指数并统一归一化, 就得到 Softmax.
        \item 因而, Softmax 可以看作 Logistic 回归在多分类情形下的直接推广.
    \end{itemize}
    \item \textbf{与交叉熵的配合}
    \begin{itemize}
        \item 设原始得分为 $\boldsymbol{z} = (z_1, \dots, z_C)$, 预测概率为
        \begin{equation*}
            \hat{y}_k = \frac{\exp(z_k)}{\sum_{j=1}^C \exp(z_j)}, \qquad k = 1, \dots, C.
        \end{equation*}
        对应的交叉熵损失为
        \begin{equation*}
            \mathcal{L} = -\sum_{k=1}^C y_k \log \hat{y}_k.
        \end{equation*}
        \item 由
        \begin{equation*}
            \log \hat{y}_k = z_k - \log \sum_{j=1}^C \exp(z_j),
        \end{equation*}
        可得
        \begin{equation*}
            \frac{\partial \mathcal{L}}{\partial z_m} = \hat{y}_m - y_m.
        \end{equation*}
        这一结果说明梯度恰好等于“预测概率减去真实概率”, 因而推导和实现都非常直接.
    \end{itemize}
    \item \textbf{物理视角}
    \begin{itemize}
        \item Softmax 的形式与统计物理中的玻尔兹曼分布一致:
        \begin{equation*}
            P(E_i) = \frac{\exp\left(-\frac{E_i}{kT}\right)}{\sum_j \exp\left(-\frac{E_j}{kT}\right)}.
        \end{equation*}
        \item 若把类别得分看成负能量, 或把温度参数吸收到得分缩放中, 就得到与 Softmax 相同的形式. 这说明 Softmax 也可以理解为一种“得分越高, 概率越大”的能量模型分布.
    \end{itemize}
    \item \textbf{历史补充}
    \begin{itemize}
        \item Softmax 这一名称通常追溯到 John S. Bridle 在 1989 年关于分类网络输出概率解释的工作\citep{bridle1990probabilistic}. 该工作明确将归一化指数映射用于多分类神经网络输出层, 也促成了 Softmax 在机器学习中的广泛使用.
    \end{itemize}
\end{enumerate}

% ===================== 感知器(Perceptron) =====================
\section{感知器(Perceptron)}
\begin{enumerate}
    \item \textbf{模型定义}
    \begin{itemize}
        \item 感知器提出于 1958 年, 是早期机器学习中的经典线性分类模型.
        \item 对二分类标签 $y \in \{+1, -1\}$, 感知器使用线性判别函数并通过符号函数作出预测:
        \begin{equation*}
            \hat{y} = \operatorname{sign}(\boldsymbol{w}^\top \boldsymbol{x}).
        \end{equation*}
    \end{itemize}
    \item \textbf{错误驱动的更新规则}
    \begin{itemize}
        \item 感知器的训练目标是找到一组参数, 使所有样本满足
        \begin{equation*}
            y^{(n)}\bigl(\boldsymbol{w}^\top \boldsymbol{x}^{(n)}\bigr) > 0.
        \end{equation*}
        \item 当样本 $(\boldsymbol{x}, y)$ 被分错, 即
        \begin{equation*}
            y\bigl(\boldsymbol{w}^\top \boldsymbol{x}\bigr) \le 0,
        \end{equation*}
        就执行更新
        \begin{equation*}
            \boldsymbol{w}_{t+1} = \boldsymbol{w}_t + y\boldsymbol{x}.
        \end{equation*}
        \item 这一更新可看作对感知器损失
        \begin{equation*}
            L_{\mathrm{perc}}(\boldsymbol{w}; \boldsymbol{x}, y) = \max\bigl(0, -y\boldsymbol{w}^\top \boldsymbol{x}\bigr)
        \end{equation*}
        的一次随机次梯度下降. 这里使用的是感知器损失, 而不是带单位间隔的 Hinge 损失.
    \end{itemize}
    \item \textbf{为何只对错分样本更新}
    \begin{itemize}
        \item 若样本已经满足 $y\boldsymbol{w}^\top \boldsymbol{x} > 0$, 则其感知器损失为 $0$, 不会推动参数继续变化.
        \item 若样本被分错, 则更新后有
        \begin{equation*}
            y\bigl(\boldsymbol{w}_{t+1}^\top \boldsymbol{x}\bigr)
            = y\bigl(\boldsymbol{w}_t^\top \boldsymbol{x}\bigr) + y^2 \boldsymbol{x}^\top \boldsymbol{x},
        \end{equation*}
        其中第二项恒非负, 因而该样本的分类裕度会朝正确方向增大.
    \end{itemize}
    \item \textbf{收敛性与局限}
    \begin{itemize}
        \item 若训练数据线性可分, 感知器算法会在有限步内收敛到某个可分解.
        \item 若数据本身不可线性分, 标准感知器通常不会收敛, 参数可能持续振荡.
    \end{itemize}
    \item \textbf{与 Logistic 回归的比较}
    \begin{itemize}
        \item 感知器只关心“是否分对”, 一旦样本处在正确一侧就停止更新.
        \item Logistic 回归则基于概率残差更新参数, 即使样本已经分对, 只要预测概率与真实标签仍有偏差, 参数就可能继续调整.
    \end{itemize}
\end{enumerate}

\textbf{问题思考}:
\begin{itemize}
    \item 感知器算法为何只对分错样本更新?
    \item 感知器收敛性的前提条件是什么?
\end{itemize}

% ===================== 支持向量机(SVM) =====================
\section{支持向量机(SVM)}
\begin{enumerate}
    \item \textbf{最大间隔思想}
    \begin{itemize}
        \item 虽然许多线性分类器都能找到把两类样本分开的超平面, 但不同超平面的鲁棒性并不相同.
        \item 支持向量机的核心思想是在所有可分超平面中寻找间隔最大的那个, 以提升模型对扰动的稳定性.
        \item 与最优超平面距离最近的样本称为\textbf{支持向量}, 它们决定了最终的决策边界.
    \end{itemize}
    \item \textbf{硬间隔 SVM}
    \begin{itemize}
        \item 设分类超平面为 $\boldsymbol{w}^\top \boldsymbol{x} + b = 0$, 对样本 $(\boldsymbol{x}^{(n)}, y^{(n)})$ 其中 $y^{(n)} \in \{+1, -1\}$, 其几何间隔可写为
        \begin{equation*}
            \gamma^{(n)} = \frac{y^{(n)}\bigl(\boldsymbol{w}^\top \boldsymbol{x}^{(n)} + b\bigr)}{\|\boldsymbol{w}\|}.
        \end{equation*}
        \item 若数据严格线性可分, 可要求所有样本满足 $y^{(n)}(\boldsymbol{w}^\top \boldsymbol{x}^{(n)} + b) \ge 1$, 并求解
        \begin{equation*}
            \begin{array}{cl}
                \min_{\boldsymbol{w}, b} & \frac{1}{2}\|\boldsymbol{w}\|^2 \\
                \text{s.t.} & y^{(n)}\bigl(\boldsymbol{w}^\top \boldsymbol{x}^{(n)} + b\bigr) \ge 1, \quad n = 1, \dots, N.
            \end{array}
        \end{equation*}
        \item 这一形式通过固定函数间隔为 $1$ 消除了缩放不确定性, 等价于最大化最小几何间隔.
    \end{itemize}
    \item \textbf{软间隔 SVM}
    \begin{itemize}
        \item 实际数据往往含有噪声, 甚至本身不可线性分. 此时可以引入松弛变量 $\xi_n \ge 0$, 允许部分样本落入间隔内部或被错分.
        \item 软间隔 SVM 的优化问题为
        \begin{equation*}
            \begin{array}{ll}
                \min_{\boldsymbol{w}, b, \boldsymbol{\xi}} & \frac{1}{2}\|\boldsymbol{w}\|^2 + C \sum_{n=1}^N \xi_n \\
                \text{s.t.} & y^{(n)}\bigl(\boldsymbol{w}^\top \boldsymbol{x}^{(n)} + b\bigr) \ge 1 - \xi_n, \quad n = 1, \dots, N, \\
                & \xi_n \ge 0, \quad n = 1, \dots, N.
            \end{array}
        \end{equation*}
        \item 当 $0 < \xi_n < 1$ 时, 样本虽然仍被分对, 但已经落入间隔内部; 当 $\xi_n \ge 1$ 时, 样本可能被错分.
        \item 超参数 $C > 0$ 用于平衡“最大化间隔”和“惩罚违约”这两个目标: $C$ 越大, 对违约样本惩罚越强; $C$ 越小, 模型越强调间隔和鲁棒性.
    \end{itemize}
    \item \textbf{Hinge 损失视角}
    \begin{itemize}
        \item 在适当重参数化下, 软间隔 SVM 等价于最小化 Hinge 损失与正则化项之和:
        \begin{equation*}
            \min_{\boldsymbol{w}, b} \sum_{n=1}^N \max\Bigl(0, 1 - y^{(n)}\bigl(\boldsymbol{w}^\top \boldsymbol{x}^{(n)} + b\bigr)\Bigr)
            + \frac{\lambda}{2}\|\boldsymbol{w}\|^2.
        \end{equation*}
        \item Hinge 损失只要分类裕度未超过 $1$ 就继续惩罚, 因而不仅要求“分对”, 还要求“分得足够远”.
    \end{itemize}
    \item \textbf{与其他分类损失的比较}
    \begin{itemize}
        \item 若记 $m = y f(\boldsymbol{x})$ 为分类裕度, 则不同损失函数对大裕度样本的处理方式不同.
        \item \textbf{平方损失}: $\mathcal{L} = (1 - m)^2$. 它对已经正确且裕度很大的样本仍会持续惩罚, 因而通常不如专门的分类损失自然.
        \item \textbf{Logistic 损失}: $\mathcal{L} = \log(1 + \exp(-m))$. 该损失光滑、便于优化, 且具有概率解释.
        \item \textbf{感知器损失}: $\mathcal{L} = \max(0, -m)$. 只要样本分对, 损失就为零.
        \item \textbf{Hinge 损失}: $\mathcal{L} = \max(0, 1 - m)$. 它比感知器损失更强调安全间隔, 因而更有利于提升鲁棒性.
    \end{itemize}
    \item \textbf{优化方法}
    \begin{itemize}
        \item 经典 SVM 常通过对偶问题转化为凸二次规划来求解, 代表算法包括 SMO(Sequential Minimal Optimization).
        \item 在大规模问题中, 也可以采用随机次梯度下降等近似优化方法.
    \end{itemize}
\end{enumerate}

\textbf{问题思考}:
\begin{itemize}
    \item SVM 为何要最大化间隔?
    \item 软间隔 SVM 如何处理噪声数据?
\end{itemize}

% ===================== 模型对比总结 =====================
\section{模型对比总结}
\begin{enumerate}
    \item \textbf{模型选择的主要考虑}
    \begin{itemize}
        \item 若需要直接输出二分类概率, 通常选择 Logistic 回归.
        \item 若需要多分类概率输出, 通常选择 Softmax 回归.
        \item 若数据线性可分且关注在线更新, 感知器是最简单的基线方法.
        \item 若更关注分类间隔与鲁棒性, 则 SVM 往往更有优势.
    \end{itemize}
    \item \textbf{典型模型对比}
    \begin{center}
    \begin{tabular}{|l|l|l|l|l|}
    \hline
    模型 & 激活函数 & 损失函数 & 优化方法 & 主要应用场景 \\
    \hline
    线性回归 & - & 均方误差 & 最小二乘/梯度下降 & 回归任务 \\
    Logistic 回归 & Sigmoid & 二元交叉熵 & 梯度下降 & 二分类概率预测 \\
    Softmax 回归 & Softmax & 多元交叉熵 & 梯度下降 & 多分类概率预测 \\
    感知器 & 符号函数(sgn) & 感知器损失 & 在线更新/随机梯度法 & 线性可分分类 \\
    SVM & 符号函数(sgn) & Hinge 损失 + $L_2$ 正则 & 二次规划/SMO/次梯度法 & 最大间隔分类 \\
    \hline
    \end{tabular}
    \end{center}
\end{enumerate}

\textbf{问题思考}:
\begin{itemize}
    \item 不同模型的损失函数有何异同?
    \item 线性模型在实际应用中应如何选择?
\end{itemize}

(注: 详细定义和更多内容请参考课程教材第 3 章及附录 E.)

\bibliographystyle{plainnat}
\bibliography{references}

\end{document}